\documentclass[12pt,a4paper]{article}

\usepackage[utf8]{inputenc}
\usepackage[T1]{fontenc}
\usepackage{amsmath,amssymb,amsthm}
\usepackage{doi}
\usepackage{hyperref}


\newtheorem{theorem}{Theorem}[section]
\newtheorem{lemma}[theorem]{Lemma}
\newtheorem{corollary}[theorem]{Corollary}
\theoremstyle{remark}
\newtheorem{remark}[theorem]{Remark}
\theoremstyle{definition}
\newtheorem{definition}[theorem]{Definition}

\DeclareMathOperator{\tr}{tr}
\newcommand{\Sym}{\mathrm{Sym}}
\newcommand{\SU}{\mathrm{SU}}

\title{Local Spectral Closure of the Capacity--Metric Bridge\\[4pt]
\normalsize\itshape LC-O2-O1 of the Lorentz--Capacity Sub-Programme}
\author{J\'er\^ome Beau\\[2pt]
\small Independent Researcher\\
\small \texttt{jerome.beau@cosmochrony.org}}
\date{2026}

\begin{document}

  \maketitle

  \begin{abstract}
    We close the local bridge between the capacity radius and the temporal co-metric
    component in the Lorentz-capacity formulation of the Cosmochrony programme.
    The local theory introduces an occupancy field $\epsilon(x)$ through the residual
    Born--Infeld capacity radius $R(x)^{2}=1-\epsilon(x)$, while the metric formulation
    introduces the same radius through $R(x)^{2}=-2g^{\tau\tau}(x)$.
    Their equivalence was previously stated as an open analytic bridge in~\cite{Beau2026lco2}.
    We prove that the identity
    \[
      \epsilon(x)=1+2g^{\tau\tau}(x)
    \]
    follows directly from local admissibility data.
    The key input is that $\epsilon(x)$ is a quadratic occupancy in the rank-three sector
    $\Sym^{2}(V_{\rho})$, normalised by the maximal Born--Infeld capacity
    $\sigma_{\mathrm{max}}=c_{\mathrm{BI}}^{2}$.
    A local version of the admissibility-weight argument of~\cite{Beau2026w1} then gives
    \[
      A_{\tau}(x)=2\bigl(1-\epsilon(x)\bigr).
    \]
    Using the principal-symbol identification of the effective operator~\cite{Beau2026q6b,Beau2026q11},
    together with the flat-space normalisation $A_{\tau}^{\infty}=2$,
    $g^{\tau\tau,\infty}=-\tfrac{1}{2}$, yields
    \[
      g^{\tau\tau}(x)=-\frac{A_{\tau}(x)}{4}.
    \]
    Therefore $R(x)^{2}=1-\epsilon(x)=-2g^{\tau\tau}(x)$.
    No use is made of the Einstein equation.
    The bridge is closed spectrally from the local admissibility form alone.
  \end{abstract}

  \medskip
  \noindent\textbf{Keywords.}
  local capacity radius; spectral occupancy; admissibility-weight reduction;
  temporal co-metric; Born--Infeld capacity; effective Lorentzian geometry;
  gravitational time dilation; Cosmochrony.

  \newpage
  \tableofcontents

  \newpage
  \section{Introduction}
    \label{sec:introduction}

    The Lorentz-capacity formulation~\cite{Beau2026lco2} introduces a local occupancy
    field $\epsilon(x)$ measuring how much of the projective capacity at $x$ is already
    locked by admissible spectral structure.
    On the capacity side, the residual radius is defined by
    \[
      R(x)^{2} := 1-\epsilon(x).
    \]
    On the metric side, the same quantity is defined by the temporal co-metric component,
    \[
      R(x)^{2} := -2g^{\tau\tau}(x).
    \]
    The open problem LC-O2-O1 of~\cite{Beau2026lco2} (Remark~3.2 and Section~7.3) is to prove
    \[
      1-\epsilon(x)=-2g^{\tau\tau}(x)
    \]
    from admissibility data alone, without importing the Einstein equation from the
    spectral gravity branch.

    The natural object controlling this bridge is the local temporal coefficient
    $A_{\tau}(x)$ of the effective admissibility operator $L_{\mathrm{eff}}$.
    In the homogeneous flat limit, this coefficient is fixed by~\cite{Beau2026q11}:
    \[
      A_{\tau}^{\infty}=2.
    \]
    When a fraction $\epsilon(x)$ of the local admissible fibre is already occupied,
    the coefficient is reduced by the free quadratic capacity fraction:
    \[
      A_{\tau}(x)=2\bigl(1-\epsilon(x)\bigr).
    \]
    This is the local counterpart of the admissibility-weight convergence argument
    of~\cite{Beau2026w1}.
    The reduction is exact only if $\epsilon(x)$ is defined as a quadratic occupancy
    in the rank-three sector $\Sym^{2}(V_{\rho})$.
    It does not hold for an arbitrary additive scalar density.

    Once this local weight reduction is established, the metric bridge follows from the
    principal-symbol reconstruction of $L_{\mathrm{eff}}$:
    \[
      g^{\tau\tau}(x)=-\frac{A_{\tau}(x)}{4}.
    \]
    The factor $\tfrac{1}{4}$ is fixed by the LCII normalisation $R^{2}=-2g^{\tau\tau}$,
    together with the flat-space values $A_{\tau}^{\infty}=2$ and
    $g^{\tau\tau,\infty}=-\tfrac{1}{2}$.
    This identification belongs to the operator-symbol reconstruction of the effective
    metric~\cite{Beau2026q6b,Beau2026q11} and does not use the Einstein equation.

    Section~\ref{sec:quadratic-occupancy} defines the quadratic local occupancy precisely.
    Section~\ref{sec:weight-reduction} proves the local weight reduction (Lemma).
    Section~\ref{sec:principal-symbol} records the principal-symbol identification.
    Section~\ref{sec:spectral-closure} states and proves the main result (Theorem and Corollary).
    Section~\ref{sec:discussion} discusses the role of the quadratic definition and the
    scope of the result.

  \section{Quadratic local occupancy}
    \label{sec:quadratic-occupancy}

    Let $V_{\rho}$ be the Weil module of $\mathrm{Heis}_{3}(\mathbb{Z}/q\mathbb{Z})$ forced
    by admissibility~\cite{Beau2026q6b}.
    The local rank-three effective sector is identified as
    \[
      H_{\mathrm{eff}}(x)\simeq \Sym^{2}(V_{\rho}),
    \]
    equipped with its $\SU(2)$-invariant quadratic norm.
    The local Born--Infeld capacity is measured in this norm.

    \begin{definition}[Local quadratic occupancy]
      \label{def:local-occupancy}
      Let $C_{\mathrm{occ}}(x)$ denote the occupied part of the local admissible sector over
      $x$, and let $C_{\mathrm{max}}(x)$ denote a maximally saturated admissible vector, with
      \[
        \sigma_{\mathrm{max}}
        :=\|C_{\mathrm{max}}(x)\|_{\Sym^{2}(V_{\rho})}^{2}
        =c_{\mathrm{BI}}^{2}.
      \]
      The local occupancy field is
      \[
        \epsilon(x)
        :=\frac{\|C_{\mathrm{occ}}(x)\|_{\Sym^{2}(V_{\rho})}^{2}}{\sigma_{\mathrm{max}}},
        \qquad 0\leq\epsilon(x)\leq 1.
      \]
      Equivalently, in any admissible orthonormal frame of $\Sym^{2}(V_{\rho})$,
      \[
        \epsilon(x)
        =\frac{\sum_{a=1}^{3}|c_{a}^{\mathrm{occ}}(x)|^{2}}
        {\sum_{a=1}^{3}|c_{a}^{\mathrm{max}}(x)|^{2}}.
      \]
    \end{definition}

    This expression is frame-independent because the norm is the $\SU(2)$-invariant
    quadratic norm on $\Sym^{2}(V_{\rho})$.
    The field $\epsilon(x)$ is not introduced as an independent scalar density.
    It is the fraction of the local admissible fibre already occupied in the quadratic norm
    that defines the Born--Infeld capacity.

  \section{Local admissibility-weight reduction}
    \label{sec:weight-reduction}

    \begin{lemma}[Local admissibility-weight reduction]
      \label{lem:local-admissibility-weight-reduction}
      Let $A_{\tau}^{\infty}=2$ be the homogeneous temporal admissibility coefficient.
      If a fraction $\epsilon(x)$ of the local $\Sym^{2}(V_{\rho})$ capacity is occupied,
      then the temporal coefficient of the local effective admissibility operator satisfies
      \[
        A_{\tau}(x)
        =A_{\tau}^{\infty}\bigl(1-\epsilon(x)\bigr)
        =2\bigl(1-\epsilon(x)\bigr).
      \]
    \end{lemma}

    \begin{proof}
      The effective operator $L_{\mathrm{eff}}$ is local at the level of its principal symbol.
      At each point $x$, its temporal coefficient is the normalised trace of the admissibility
      weight operator $\mathrm{Ad}$ over the fibre sector available over $x$~\cite{Beau2026w1}.
      In the homogeneous case, the full image sector is available, and the fibre-averaging
      argument of~\cite{Beau2026w1} gives $A_{\tau}^{\infty}=2$.

      Locally, the occupied and free parts of the fibre decompose orthogonally in the
      $\SU(2)$-invariant quadratic norm on $\Sym^{2}(V_{\rho})$:
      \[
        C_{\mathrm{max}}(x)=C_{\mathrm{occ}}(x)\oplus C_{\mathrm{free}}(x).
      \]
      Therefore $\|C_{\mathrm{free}}(x)\|^{2}=\|C_{\mathrm{max}}(x)\|^{2}-\|C_{\mathrm{occ}}(x)\|^{2}$,
      and dividing by $\sigma_{\mathrm{max}}$ gives
      \[
        \frac{\|C_{\mathrm{free}}(x)\|^{2}}{\|C_{\mathrm{max}}(x)\|^{2}}=1-\epsilon(x).
      \]

      Since $\Sym^{2}(V_{\rho})$ is an irreducible $\SU(2)$-representation and the
      admissibility weight $\mathrm{Ad}$ is $\SU(2)$-equivariant, Schur's lemma forces
      $\mathrm{Ad}=\lambda\cdot\mathrm{Id}$ on this sector for some scalar $\lambda>0$.
      Therefore the partial trace over any subspace $C\subseteq\Sym^{2}(V_{\rho})$ satisfies
      \[
        \tr_{C}(\mathrm{Ad})=\lambda\|C\|^{2}.
      \]
      Consequently,
      \[
        \frac{\tr_{C_{\mathrm{free}}(x)}(\mathrm{Ad})}{\tr_{C_{\mathrm{max}}(x)}(\mathrm{Ad})}
        =\frac{\|C_{\mathrm{free}}(x)\|^{2}}{\|C_{\mathrm{max}}(x)\|^{2}}
        =1-\epsilon(x).
      \]
      The local temporal coefficient is the normalised trace over $C_{\mathrm{free}}(x)$:
      \[
        A_{\tau}(x)
        =A_{\tau}^{\infty}
        \frac{\tr_{C_{\mathrm{free}}(x)}(\mathrm{Ad})}{\tr_{C_{\mathrm{max}}(x)}(\mathrm{Ad})}
        =A_{\tau}^{\infty}\bigl(1-\epsilon(x)\bigr).
      \]
      Using $A_{\tau}^{\infty}=2$ yields $A_{\tau}(x)=2\bigl(1-\epsilon(x)\bigr)$.
    \end{proof}

  \section{Principal-symbol identification}
    \label{sec:principal-symbol}

    The effective metric is reconstructed from the principal symbol of $L_{\mathrm{eff}}$~\cite{Beau2026q6b}.
    For the temporal component, in the LCII normalisation~\cite{Beau2026lco2}, this gives
    \[
      g^{\tau\tau}(x)=-\frac{A_{\tau}(x)}{4}.
    \]
    This identification follows from the principal-symbol formula of~\cite{Beau2026q6b},
    applied locally, together with the flat-space normalisation $A_{\tau}^{\infty}=2$,
    $g^{\tau\tau,\infty}=-\tfrac{1}{2}$ established in~\cite{Beau2026q11}.

    This step is purely operator-theoretic.
    It uses the local principal symbol of the admissibility operator and does not use the
    Einstein equation, the spectral-entropy variation, or the field equations of
    the Gravity paper~\cite{Beau2026q6b}.

  \section{Spectral closure of the capacity--metric bridge}
    \label{sec:spectral-closure}

    \begin{theorem}[Spectral closure of the local capacity--metric bridge]
      \label{thm:spectral-closure}
      The local capacity radius and the metric radius coincide:
      \[
        R(x)^{2}=1-\epsilon(x)=-2g^{\tau\tau}(x).
      \]
      Equivalently,
      \[
        \epsilon(x)=1+2g^{\tau\tau}(x).
      \]
    \end{theorem}

    \begin{proof}
      By Definition~\ref{def:local-occupancy}, $R(x)^{2}=1-\epsilon(x)$.
      By Lemma~\ref{lem:local-admissibility-weight-reduction},
      $A_{\tau}(x)=2\bigl(1-\epsilon(x)\bigr)$.
      The local principal-symbol identification of Section~\ref{sec:principal-symbol} gives
      $g^{\tau\tau}(x)=-A_{\tau}(x)/4$.
      Hence
      \[
        -2g^{\tau\tau}(x)=\frac{A_{\tau}(x)}{2}=1-\epsilon(x).
      \]
      Therefore $R(x)^{2}=1-\epsilon(x)=-2g^{\tau\tau}(x)$, and rearranging gives the
      equivalent form.
    \end{proof}

    \begin{corollary}[Upgrade of Remark~3.2 of \cite{Beau2026lco2}]
      \label{cor:lcii-upgrade}
      The equivalence between the capacity and metric definitions of $R(x)$,
      postulated in Remark~3.2 of~\cite{Beau2026lco2}, is not an independent postulate.
      It follows from local admissibility-weight reduction on the quadratic fibre sector
      $\Sym^{2}(V_{\rho})$.
      The chain
      \[
        \epsilon(x)
        \longmapsto A_{\tau}(x)
        \longmapsto g^{\tau\tau}(x)
      \]
      is closed spectrally, without importing the Einstein equation.
    \end{corollary}

    \begin{remark}[Role of the quadratic occupancy]
      \label{rem:role-quadratic-occupancy}
      The proof depends essentially on $\epsilon(x)$ being a quadratic occupancy in
      $\Sym^{2}(V_{\rho})$.
      If $\epsilon(x)$ were only an additive scalar density, the Schur step would not apply,
      the orthogonal decomposition
      $C_{\mathrm{max}}(x)=C_{\mathrm{occ}}(x)\oplus C_{\mathrm{free}}(x)$
      would not imply the exact factor $1-\epsilon(x)$ in $A_{\tau}(x)$, and the bridge
      would remain open.
      The identity $\epsilon(x)=1+2g^{\tau\tau}(x)$ is therefore a spectral admissibility
      identity, not a phenomenological identification between matter density and metric
      curvature.
    \end{remark}

  \section{Discussion}
    \label{sec:discussion}

    The result shows that $\epsilon(x)$ and $g^{\tau\tau}(x)$ are not independent fields:
    they are two readings of the same local admissibility datum.
    The projective capacity available at $x$ for further temporal transitions is
    $1-\epsilon(x)$, and this is precisely the quantity that the effective metric encodes
    in its temporal component $g^{\tau\tau}(x)=-\tfrac{1}{2}(1-\epsilon(x))$.

    Gravitational time dilation then takes the following form in the capacity language.
    A massive configuration at $x$ occupies a fraction $\epsilon(x)>0$ of the local BI
    budget; the residual temporal capacity $R(x)=\sqrt{1-\epsilon(x)}<1$ slows local
    clocks relative to the vacuum.
    This is not imported from the Einstein equation.
    It is a consequence of the reduction of the available admissible fibre sector.

    The Schwarzschild case provides a consistency check via~\cite{Beau2026lco2},
    Corollary~6.1: in the weak-field regime,
    $R(x)\approx 1+\Phi_{N}(x)/c^{2}$, so $\epsilon(x)\approx -2\Phi_{N}(x)/c^{2}$, with
    $\Phi_{N}<0$ near a mass source.
    The capacity budget is slightly more occupied near a massive body, reducing the temporal
    residual, consistently with the standard gravitational redshift.

  \section{Conclusion}
    \label{sec:conclusion}

    We have closed LC-O2-O1.
    The local capacity radius $R(x)^{2}=1-\epsilon(x)$ and the metric radius
    $R(x)^{2}=-2g^{\tau\tau}(x)$ are the same quantity, once the occupancy field is
    defined as a quadratic occupancy in $\Sym^{2}(V_{\rho})$.

    The proof proceeds through two steps.
    First, the local temporal coefficient of the admissibility operator is reduced by the
    free quadratic capacity fraction:
    \[
      A_{\tau}(x)=2\bigl(1-\epsilon(x)\bigr).
    \]
    This follows from the $\SU(2)$-invariance of the admissibility weight (Schur's lemma)
    and the orthogonality of the occupied/free decomposition.
    Second, the principal-symbol identification gives
    \[
      g^{\tau\tau}(x)=-\frac{A_{\tau}(x)}{4}.
    \]
    Combining these two identities yields $R(x)^{2}=1-\epsilon(x)=-2g^{\tau\tau}(x)$.
    The bridge is closed spectrally, without the Einstein equation.
    The open bridge of Remark~3.2 of~\cite{Beau2026lco2} is hereby upgraded to a theorem.

    \newpage
    \bibliographystyle{plain}
    \bibliography{cosmochrony-bibliography}

\end{document}
